%
% FH Technikum Wien
% !TEX encoding = UTF-8 Unicode
%
% Version V2.24 von 2024-12-19 otrebski
%
\documentclass[BMR,Seminar,ngerman,IEEE]{twbook}
\renewcommand*{\degreecourse}{FH Technikum Wien, Bachelorstudiengang BIF\\Lehrveranstaltung BIF-VZ-5-WS2025-INFC-EN}

% --- Encoding & fonts ---
\usepackage[utf8]{inputenc}
\usepackage[T1]{fontenc}

% --- Graphics, figures, urls, lists, tables, code ---
\usepackage{float}
\usepackage{placeins}
\usepackage{graphicx}
\usepackage{caption}
\usepackage{booktabs}
\usepackage{tabularx}
\usepackage{enumitem}
\usepackage{listings}
\usepackage[final]{microtype} % hilft gegen Overfull/Underfull
\emergencystretch=2em         % sanfter Puffer für Zeilenumbruch

% Screenshots hier hinein legen:
\graphicspath{{screenshots/}}

% --- Bibliography (twbook lädt biblatex via Hook, sobald inputenc geladen ist) ---
\addbibresource{Literatur.bib}

% --- Listings (inkl. YAML) ---
\lstdefinelanguage{yaml}{
  keywords={true,false,null,y,n},
  comment=[l]{\#},
  morestring=[b]',
  morestring=[b]",
  sensitive=false,
  showstringspaces=false,
  literate =    {---}{{\textcolor{gray}{---}}}3
                {>}{{\textcolor{gray}{>}}}1
                {|}{{\textcolor{gray}{|}}}1
                {:}{{\textcolor{blue}{:}}}1
                {-}{{\textcolor{gray}{-}}}1,
}
\lstset{
  basicstyle=\ttfamily\small,
  frame=single,
  breaklines=true,
  backgroundcolor=\color{gray!5},
  captionpos=b
}

% --- Convenience figure macros ---
\newcommand{\screenshotH}[3]{%
  \begin{figure}[H]
    \centering
    \includegraphics[width=\linewidth]{#1}%
    \caption{#2}%
    \label{fig:#3}%
  \end{figure}%
}
\newcommand{\placeholderfig}[2]{%
  \begin{figure}[htbp]
    \centering
    \fbox{\rule{0pt}{120pt}\rule{0.9\linewidth}{0pt}}%
    \caption{#1}%
    \label{fig:#2}%
  \end{figure}%
}

% Die nachfolgenden Pakete stellen sonst nicht benötigte Features zur Verfügung
\usepackage{blindtext}

% --- Hyperref am Schluss laden (robuster) ---
\usepackage{hyperref}

%
% Einträge für Deckblatt, Kurzfassung, etc.
%
\title{Ansible Workshop Protocol}
\author{David Veigel \and Kevin Forter}
\studentnumber{XXXXXXXXXXXXXXX} % ggf. beide Nummern: 12345678, 87654321
\place{Wien}
% \acknowledgements{\blindtext}

\begin{document}

\maketitle

% ============================================================================
\onecolumn
\newpage

% ============================================================================
\chapter{Ansible Workshop}
In this workshop, we explore the capabilities of Ansible for automation and configuration management. We set up a master node and worker nodes, configure SSH access, and deploy playbooks to manage files and web servers.

\section{Installation and Setup}

\subsection{Installing Ansible}
First, we install Ansible on the master node. This involves updating the package manager and installing the ansible package.
\screenshotH{master_install_ansible.png}{Installing Ansible on Master}{install}

\subsection{SSH Key Generation}
To allow passwordless authentication between the master and worker nodes, we generate an SSH key pair on the master node.
\screenshotH{master_ssh_key_generation.png}{Generating SSH Keys}{ssh-gen}

\subsection{Distributing Keys}
The public key is then copied to the worker nodes. This ensures that the master can connect to the workers without manual password entry.
\screenshotH{worker_save_master_ssh_key.png}{Saving Master Key on Worker}{ssh-copy}

\subsection{Verifying Connection}
We verify the SSH connection to the worker nodes to ensure that the setup is correct.
\screenshotH{master_connection_worker1_succesfull.png}{Connection to Worker 1 Successful}{conn-w1}
\screenshotH{master_connection_worker2_succesfull.png}{Connection to Worker 2 Successful}{conn-w2}

\section{Configuration}

\subsection{Creating Configuration}
We create a directory for our Ansible project and generate a default configuration file.
\screenshotH{master_create_ansible_config.png}{Creating Ansible Config}{create-conf}

\subsection{Ansible Config File}
The \texttt{ansible.cfg} file is customized to define the inventory file location and other defaults.
\screenshotH{master_ansible_config.png}{Ansible Configuration File}{ansible-conf}

\subsection{Hosts Inventory}
We define our inventory in the \texttt{hosts} file, listing the IP addresses of our worker nodes and grouping them if necessary.
\screenshotH{master_ansible_hosts_config.png}{Hosts Inventory Configuration}{hosts-conf}

\subsection{Ping Test}
We use the \texttt{ansible} command with the \texttt{ping} module to test connectivity to all hosts in our inventory.
\screenshotH{master_ping_workers.png}{Pinging Workers}{ping}

\section{Playbooks}

\subsection{File Operations}
We start with a simple playbook to copy a file to the worker nodes.
\screenshotH{master_create_testfile_copy_to_worker.png}{Creating Copy Playbook}{copy-playbook}

We run the playbook and observe the output to ensure the task is changed.
\screenshotH{master_check_for_testfile_copy.png}{Running Copy Playbook}{run-copy}

We verify on the worker node that the file has indeed been copied.
\screenshotH{worker_after_testfile_copy.png}{File Copied to Worker}{verify-copy}

We also create a playbook to delete the file and clean up.
\screenshotH{master_delete_testfile_on_worker.png}{Deleting File Playbook}{delete-playbook}

\subsection{Webserver Deployment}
Next, we create a more complex playbook to install and configure a web server (Apache/Nginx) on the workers.
\screenshotH{master_webserver_config_yaml.png}{Webserver Configuration Playbook}{web-playbook}

We first deploy this playbook to a single worker to test.
\screenshotH{master_load_playbook_on_worker1.png}{Deploying to Worker 1}{deploy-w1}

Once satisfied, we deploy it to all worker nodes.
\screenshotH{master_load_playbook_on_all_worker.png}{Deploying to All Workers}{deploy-all}

\subsection{Verification}
Finally, we verify the deployment by accessing the web server from a browser.
\screenshotH{browser_worker1_website_success.png}{Website Successfully Deployed}{web-success}

\end{document}
